\section{Methodology}
\label{sec:methodology}

\subsection{The \gore{} Language}
\label{sec:language}

\gore{} is a declarative, relation-based language in the logic-programming
tradition, distinguished by making the search tree a first-class syntactic
object, augmented with a functional sublanguage for deterministic computation
and explicit effect annotations for non-deterministic external calls.

\paragraph{Grammar.}
\todo{Insert BNF grammar table: Program, Clause, Body, Primitive, Term.}

\paragraph{Five primitives.}
\begin{description}
  \item[\textsc{Step}$(t_1, t_2)$] Unification / result binding.
    Succeeds if $t_1$ and $t_2$ unify under the current substitution; extends
    the substitution on success and fails otherwise.
    Corresponds to deterministic forward progress; no choice point is created.
  \item[\textsc{Fork}$(B_1, \ldots, B_k)$] Non-deterministic branching.
    Creates $k$ choice points and evaluates each branch $B_i$ in order
    (depth-first, left-to-right by default).
    On failure of branch $B_i$ the interpreter backtracks and tries $B_{i+1}$.
  \item[\textsc{Cut}$(B, \ell)$] Labelled branch pruning.
    Evaluates body $B$; if $B$ fails, jumps to label $\ell$, pruning all
    open choice points between the current point and $\ell$.
    Semantics: deterministic abort, not backtracking.
  \item[\letstmt{} $x \mathrel{:=} e$] Deterministic functional computation.
    Evaluates expression $e$ (arithmetic, string, list operations) and binds
    the result to $x$ with no backtrack point.
  \item[\callstmt{} $r \mathrel{=} @f(a_1, \ldots, a_n)$] External effect.
    Invokes external function $f$ with arguments $a_i$, binds the result to $r$,
    and logs the call as a $(f, \mathbf{a}, r)$ node in the execution trace.
    In synthetic training, $f$ is a mock oracle; at inference time, $f$ can be a
    tool, API, or sub-agent.
\end{description}

\paragraph{Operational semantics.}
\todo{Insert small-step inference rules using \texttt{mathpartir} or
      \texttt{proof} environment.  One rule per primitive.}

\subsection{Interpreter as Process Reward Model}
\label{sec:interpreter-prm}

The \gore{} interpreter evaluates a predicted trace node-by-node against the
gold trace and returns a dense scalar reward at each step.
No human annotation and no separately trained reward model are required.

\begin{verbatim}
def gore_reward(predicted_trace, correct_trace):
    step_rewards = []
    for pred, gold in zip(predicted_trace, correct_trace):
        if pred.type == gold.type:       # correct node type
            r = 0.3
            if pred.env == gold.env:     # correct environment binding
                r = 0.7
            if pred == gold:             # exact match
                r = 1.0
        else:
            r = -0.5
        step_rewards.append(r)
    return mean(step_rewards)
\end{verbatim}

This function is the PRM.
The three reward levels (0.3 / 0.7 / 1.0) provide a curriculum within a single
trace: the model receives credit for getting the node type right, additional
credit for correct environment state, and full credit for exact prediction.
Incorrect node types are penalised at $-0.5$.

\todo{Formalise as a definition; add discussion of why this decomposition
      (type $\to$ env $\to$ exact) captures the relevant structure.}

\subsection{Structural Isomorphism with Lean/Coq Tactics}
\label{sec:isomorphism}

The five \gore{} primitives are structurally isomorphic to five standard
Lean\,/\,Coq tactic combinators.
This is not a metaphor but a design decision: any model pretrained on tactic
traces already has an inductive bias toward tree-structured search that
transfers to \gore{}.

\begin{table}[h]
  \centering
  \caption{Structural isomorphism between \gore{} primitives and Lean\,/\,Coq
           tactics (left), and between \gore{} primitives and the Decomposer
           framework (right).}
  \label{tab:isomorphism}
  \begin{tabular}{lll}
    \toprule
    Lean\,/\,Coq tactic & \gore{} primitive & Decomposer concept \\
    \midrule
    \texttt{exact e}       & \textsc{Step}  & \texttt{return result} \\
    \texttt{cases h}       & \textsc{Fork}  & \texttt{raise NotImplementedError} \\
    \texttt{contradiction} & \textsc{Cut}   & \texttt{refuse} \\
    \texttt{have h := e}   & \letstmt{}     & \texttt{x = expr} \\
    \texttt{apply f}       & \callstmt{}    & \texttt{tool\_call(args)} \\
    \bottomrule
  \end{tabular}
\end{table}

\todo{Add formal definition of the isomorphism as a functor between the
      respective operational semantics.}

\subsection{Training Pipeline}
\label{sec:training}

We use a three-stage pipeline.

\paragraph{Stage 1: Formal proof pretraining.}
We fine-tune on tactic traces extracted from LeanDojo
\citep{yang2023leandojo} (${\approx}98$K theorems) and CoqGym
\citep{yang2019coqgym} (${\approx}71$K proofs), using the isomorphism in
Table~\ref{tab:isomorphism} to convert each tactic to its \gore{} counterpart
via \texttt{lean2gore.py} and \texttt{coqgym2gore.py}.
The goal is to instil an inductive bias for tree-structured search before the
model sees any \gore{} syntax.

\paragraph{Stage 2: \gore{} curriculum fine-tuning.}
We fine-tune on synthetic \gore{} programs generated by \texttt{goregen.py} at
five difficulty levels (Section~\ref{sec:curriculum}).
The model learns to predict full execution traces, supervised by maximum
likelihood.

\paragraph{Stage 3: Dr.\,GRPO with dense interpreter reward.}
We continue training with Dr.\,GRPO using the dense reward function from
Section~\ref{sec:interpreter-prm}.
Dr.\,GRPO corrects the selection bias of standard GRPO that arises because
deeper \gore{} programs admit exponentially more valid traces.
We ablate with GRPO + sparse reward (E2a) and DAPO (E2c).

\todo{Add pseudocode for the full training loop; clarify model size and
      hardware budget.}

\subsection{Curriculum Design}
\label{sec:curriculum}

The curriculum is parameterised by depth $d \in \{1,2,3,4,5\}$ and branching
factor $w \in \{1,2,3\}$.
A level-$\ell$ program has $d \leq \ell + 1$ and $w \leq \ell + 1$, giving
five difficulty levels (0--4).
Task families: enumeration (level 0--1), graph reachability (level 1--2), list
membership (level 2--3), recursive decomposition (level 3--4).

\todo{Add table: level $\times$ (d, w, \# atoms, expected backtrack depth,
      sample count).}
